\section{AGENTS.md：持久化指导}
%% ============================================================

\subsection{什么是 AGENTS.md}

\texttt{AGENTS.md} 是一个开放格式的文件，可以理解为\textbf{给 agent 看的 README}。它会被自动加载到上下文中，是编码团队标准化 Codex 行为的最佳位置。

\begin{importantbox}{AGENTS.md 应覆盖的内容}
\begin{itemize}[nosep]
    \item 仓库布局和重要目录
    \item 项目运行方式
    \item Build、test、lint 命令
    \item 工程规范和 PR 期望
    \item 约束和"不要做"规则
    \item "完成"的定义及验证方式
\end{itemize}
\end{importantbox}

\subsection{多级 AGENTS.md 体系}

Codex 支持在多个层级放置 \texttt{AGENTS.md}，形成层级覆盖机制：

\begin{itemize}
    \item \textbf{全局级}：\texttt{\textasciitilde/.codex/AGENTS.md}，存放个人默认偏好
    \item \textbf{仓库级}：项目根目录的 \texttt{AGENTS.md}，存放团队共享标准
    \item \textbf{子目录级}：特定目录中的 \texttt{AGENTS.md}，存放局部规则
\end{itemize}

\begin{knowledgebox}{就近优先原则}
当多个层级存在 \texttt{AGENTS.md} 时，离当前工作目录更近的文件优先级更高。这类似于 \texttt{.gitignore} 的覆盖逻辑。
\end{knowledgebox}

\subsection{编写原则}

\begin{itemize}
    \item \textbf{保持简洁}：短而准确的 \texttt{AGENTS.md} 比冗长模糊的规则有用得多
    \item \textbf{按需添加}：先写基础内容，只在发现重复性错误后才添加新规则
    \item \textbf{拆分引用}：当文件过大时，保持主文件精简，并引用任务级别的 markdown 文件（如 code\_review.md、architecture.md）
    \item \textbf{回顾驱动}：当 Codex 犯同样的错误两次时，让它做 retrospective 并更新 \texttt{AGENTS.md}
\end{itemize}

\begin{warningbox}{常见错误：在 Prompt 中堆积持久性规则}
把本应写在 \texttt{AGENTS.md} 中的持久性规则反复写在 prompt 里，是使用 Codex 时最常见的反模式之一。应当将经过验证的 prompting 模式迁移到 \texttt{AGENTS.md} 中。
\end{warningbox}

CLI 中可以使用 \texttt{/init} 命令快速生成一份 starter \texttt{AGENTS.md}。

\subsection{本章小结}

\texttt{AGENTS.md} 是 Codex 持久化指导的核心机制。通过多级文件体系实现全局、仓库和目录级别的规则覆盖。保持文件简洁、按需迭代，并利用 retrospective 机制持续改进。

%% ============================================================
